\chapter{System Overview and Architecture}

\section{End-to-End Workflow}

The deployed path is organized as a sequence of separable stages:

\begin{enumerate}
    \item A wearable or phone camera supplies an egocentric video stream.
    \item Frames are sampled and grouped into short overlapping clips.
    \item Two frozen video encoders convert each clip into complementary
    motion and appearance features.
    \item The fused features are processed by a causal temporal model.
    \item A fixed-lag decoder stabilizes the sequence and commits a step only
    when sufficient evidence has accumulated.
    \item The current step, completed steps, and procedure state are exposed
    through an API and rendered by the replay, live, or HoloLens client.
\end{enumerate}

\section{Streaming Model Path}

The live path samples the source at approximately 10 frames per second. Sixteen
sampled frames form a 1.6-second clip. A new clip is evaluated every 0.2
seconds, so predictions are refreshed frequently while each prediction still
contains motion context. The two feature streams are normalized and
concatenated into one 2,176-dimensional vector:

\begin{center}
\begin{tcolorbox}[colback=darkBlue!5!white,colframe=darkBlue,width=0.94\textwidth]
\centering
\textbf{10 fps frames}
\ $\rightarrow$\ 
\textbf{16-frame / 1.6 s clips}
\ $\rightarrow$\ 
\textbf{frozen motion + appearance encoders}
\ $\rightarrow$\ 
\textbf{2,176-d fusion}
\ $\rightarrow$\ 
\textbf{causal TCN}
\ $\rightarrow$\ 
\textbf{fixed-lag Viterbi}
\ $\rightarrow$\ 
\textbf{committed step}
\end{tcolorbox}
\end{center}

The motion stream is produced by VideoMAEv2-giant and has 1,408 dimensions.
The appearance stream is produced by InternVideo2-B14 and has 768 dimensions.
The encoders are frozen during the final temporal-model training; the temporal
head learns how the fused evidence evolves over the procedure. The two model
configurations are documented as dedicated views in Chapter~\ref{ch:model-architecture};
this chapter focuses on the surrounding system and application interfaces.

\section{How a New Step Is Committed}

The classifier produces a prediction at every 0.2-second update, but the
interface does not treat every changing prediction as a new step. The decoder
keeps a short history of recent class scores and applies a fixed lag of ten
updates, approximately two seconds. It selects a temporally consistent path
through that history and commits the oldest stable decision. A new step is
therefore known when the committed label changes, not merely when the
instantaneous classifier output changes.

This design reduces flicker caused by hand occlusion or ambiguous frames. The
trade-off is intentional: the demonstrator reports a step after a short
evidence-collection delay. The measured end-to-end announcement delay on the
live path is approximately 2.84 seconds.

\section{Implementation Map}

\begin{table}[H]
\centering
\small
\begin{tabularx}{\textwidth}{@{}p{0.31\textwidth}X@{}}
\toprule
\textbf{Project file or area} & \textbf{Higher-level responsibility} \\
\midrule
\texttt{copilot\_model/config.py} & Central paths, class names, feature and
checkpoint configuration. \\
\texttt{copilot\_model/encoders.py} & Loads the two pretrained encoders,
preprocesses clips, and creates the 2,176-dimensional fusion feature. \\
\texttt{copilot\_model/nets/causal\_tcn.py} & Defines the causal temporal
convolutional network used for online classification. \\
\texttt{copilot\_model/causal\_decode.py} & Runs causal scoring and fixed-lag
sequence decoding for evaluation or replay. \\
\texttt{training/train.py} and \texttt{training/augment.py} & Trains the
temporal head and creates training variations for robustness. \\
\texttt{training/eval\_gate.py} & Applies the normal, scrambled, and cold-start
acceptance checks. \\
\texttt{replay-demo/serve\_demo.py} & Provides browser replay and visualization
of the inference timeline. \\
\texttt{holo-lens-api-app/server.py} & Receives live frames, performs
inference, and publishes state to clients. \\
\bottomrule
\end{tabularx}
\caption{Functional map of the main model and application files.}
\label{tab:file-map}
\end{table}

\section{Application Services}

The project separates inference from presentation. The replay service runs on
port 8099 and allows a recorded asset to be processed and inspected. The live
service uses the camera/WebRTC path and exposes state through the API, with the
inference applications using ports 8555 and 8556 in the current setup. The
Unity HoloLens application consumes the state and renders the operator-facing
overlay; it does not contain the trained model.

\section{State Delivered to the Interfaces}

The inference service publishes a compact procedure state rather than exposing
raw model tensors to the clients. The state contains the current committed
label, the recognized procedure direction, the progress of the step list, and
status information used to distinguish normal progress from out-of-order
activity. The replay dashboard uses this state to draw the current-step panel
and the human-versus-AI timeline. The live and HoloLens clients use the same
conceptual state, which keeps presentation changes independent from model
changes.

\section{Operational Data Path}

The phone or camera is responsible for publishing the video. The server
receives frames, performs the sampling and buffering required by the model,
maintains the temporal state, and returns the latest committed result. This
separation is useful for the demonstration because the same trained checkpoint
can be tested in replay and live modes. It also makes the eventual deployment
choices explicit: inference can remain on a workstation or server while the
wearable device is used primarily for capture and display.
